% TikZ 图：节点电流守恒与回路电压守恒
% Day 19：直流电路
\begin{tikzpicture}[scale=0.92,>=Latex]
  \tikzset{
    wire/.style={thick},
    current/.style={->,very thick,red},
    volt/.style={->,thick,blue},
    res/.style={thick,fill=white},
    note/.style={font=\small}
  }

  % 左：节点电流守恒
  \begin{scope}
    \node[font=\bfseries] at (1.9,1.75) {节点电流守恒};
    \coordinate (O) at (2,0);
    \filldraw[black] (O) circle (2.2pt);

    \draw[wire] (-0.2,1.0) -- (O);
    \draw[wire] (-0.2,-1.0) -- (O);
    \draw[wire] (O) -- (4.2,0);

    \draw[current] (0.15,0.82) -- (1.25,0.27);
    \node[red,above left] at (0.50,0.88) {$I_1$};
    \draw[current] (0.15,-0.82) -- (1.25,-0.27);
    \node[red,below left] at (0.50,-0.88) {$I_2$};
    \draw[current] (2.45,0.22) -- (3.55,0.22);
    \node[red,above] at (3.00,0.24) {$I_3$};

    \node[note,align=center] at (2,-1.55) {节点不积累电荷\\[2pt]$I_1+I_2=I_3$};
  \end{scope}

  % 右：回路电压守恒
  \begin{scope}[xshift=5.6cm]
    \node[font=\bfseries] at (2.45,2.05) {回路电压守恒};

    \draw[wire] (0,-0.58) -- (0,1.15) -- (1.0,1.15);
    \draw[res] (1.0,0.85) rectangle (2.0,1.45);
    \node at (1.5,1.15) {$R_1$};
    \draw[wire] (2.0,1.15) -- (3.1,1.15);
    \draw[res] (3.1,0.85) rectangle (4.1,1.45);
    \node at (3.6,1.15) {$R_2$};
    \draw[wire] (4.1,1.15) -- (5.0,1.15) -- (5.0,-1.05) -- (0,-1.05) -- (0,-0.82);

    % 电源
    \draw[thick] (-0.45,-0.58) -- (0.45,-0.58);
    \draw[thick] (-0.28,-0.82) -- (0.28,-0.82);
    \node[left] at (-0.55,-0.70) {$U_{\mathrm s}$};
    \node[right] at (0.48,-0.58) {$+$};
    \node[right] at (0.31,-0.82) {$-$};
    \node[note,anchor=west] at (0.92,-0.70) {跨过电源：电势升高};

    \draw[current] (0.38,1.58) -- (0.92,1.58);
    \node[red,above] at (0.65,1.58) {$I$};
    \draw[volt] (1.0,0.63) -- (2.0,0.63);
    \node[blue,below] at (1.5,0.58) {$U_1$};
    \draw[volt] (3.1,0.63) -- (4.1,0.63);
    \node[blue,below] at (3.6,0.58) {$U_2$};

    \node[note,align=center] at (2.55,-1.70) {沿电流绕行一周\\[2pt]$U_{\mathrm s}=U_1+U_2$};
  \end{scope}
\end{tikzpicture}
